\documentclass[11pt,oneside]{article}
\usepackage[margin=1in]{geometry}
\usepackage{amsmath,amssymb,mathtools}
\usepackage{xcolor}
\usepackage{listings}
\usepackage{fancyhdr}
\usepackage{hyperref}

\lstset{
  basicstyle=\ttfamily\small,
  breaklines=true,
  columns=fullflexible,
  keywordstyle=\color{blue},
  commentstyle=\color{gray},
  stringstyle=\color{red},
  backgroundcolor=\color{lightgray!10}
}

\pagestyle{fancy}
\fancyhead[R]{Doc\_1: Predicate Decomposition \& Formal Spec}
\fancyfoot[C]{\thepage}

\title{Equal-Status Memory Controller: Predicate Decomposition \& Formal Specification\\
\large Field Classifier, Pair Matcher, Discrimination Detector, Registry Controller}
\author{Universal Standard Axiom Corporation}
\date{October 5, 2026}

\begin{document}

\maketitle

\begin{abstract}
This document decomposes the master safety invariant $I_0$ from Doc\_0 into four functional predicates, each assigned to a dedicated hardware sub-unit. Each predicate is formally specified using SystemVerilog Assertions (SVA), temporal logic, and transfer-function semantics. Continuous grounding ensures every state transition preserves $I_0$.
\end{abstract}

\tableofcontents
\newpage

\section{Predicate Decomposition (from $I_0$)}

The master invariant $I_0$ is decomposed into four \textbf{necessary and sufficient} predicates:

\[
I_0 \equiv P_1 \land P_2 \land P_3 \land P_4
\]

\begin{enumerate}
\item $P_1$: \textbf{Field-Classifier Correctness} (extract perturbing signature)
\item $P_2$: \textbf{Pair-Matcher Consistency} (reuse cached decisions)
\item $P_3$: \textbf{Discrimination-Detector Refusal Tracking} (prevent selective refusal)
\item $P_4$: \textbf{Registry-Controller Monotonicity} (enforce tightening-only updates)
\end{enumerate}

Failure of any single predicate violates $I_0$. All four must be maintained throughout execution.

\section{Predicate $P_1$: Field-Classifier Correctness}

\subsection{Informal Specification}

The field-classifier sub-unit reads an incoming memory operation and extracts only the fields registered as perturbing. It computes:

\[
\text{Perturb}(O) = \bigoplus_{f_i \in R} f_i
\]

where $R$ is the current perturbing set (from registry) and $\bigoplus$ is bitwise XOR.

\subsection{Formal Specification (P1)}

\begin{boxed}
\textbf{Predicate $P_1$ (Field-Classifier Correctness):}

For every operation $O$ arriving at time $t$ with fields $O.f_1, \ldots, O.f_n$:

\[
\text{sig}_t = \bigoplus_{i=1}^{n} (O.f_i \cdot R_t[i])
\]

where $R_t[i] \in \{0, 1\}$ indicates whether field $i$ is in the perturbing set at time $t$.

The classifier output $\text{sig}_t$ is latched and becomes available at the next cycle.

\textbf{Invariant:} $\text{sig}_t$ is deterministic and independent of non-perturbing fields.
\end{boxed}

\subsection{SystemVerilog Assertions for $P_1$}

\begin{lstlisting}[language=Verilog]
// SVA1: Perturbing signature is XOR of selected fields
property P1_signature_correct;
  @(posedge clk) 
  sig_next == (op_in.requester_in & REG_perturb[0] ^
               op_in.address_in   & REG_perturb[1] ^
               op_in.origin_in    & REG_perturb[2] ^
               op_in.class_in     & REG_perturb[3] ^
               op_in.epoch_in     & REG_perturb[4]);
endproperty
assert property (P1_signature_correct);

// SVA2: Non-perturbing fields do not affect signature
property P1_non_perturb_invariant;
  @(posedge clk)
  if (REG_perturb[0] == 1'b0)  // requester not perturbing
    (sig_next ^ sig_prev_alt) == 0;  // sig same if requester differs
endproperty
assert property (P1_non_perturb_invariant);

// SVA3: Signature latched at end of cycle
property P1_latch;
  @(posedge clk) sig_stable_until_next_op;
endproperty
assert property (P1_latch);
\end{lstlisting}

\section{Predicate $P_2$: Pair-Matcher Consistency}

\subsection{Informal Specification}

The pair-matcher maintains a cache of 64 slots, indexed by hash of the perturbing signature:

\[
\text{slot} = \text{hash}(\text{Perturb}(O)) \mod 64
\]

For each operation:
\begin{itemize}
\item Compute the signature (from P1).
\item Hash the signature.
\item Look up the cache.
\item If found and signature matches exactly, reuse the cached decision.
\item If not found or mismatch, forward the operation to the discrimination-detector.
\end{itemize}

\subsection{Formal Specification (P2)}

\begin{boxed}
\textbf{Predicate $P_2$ (Pair-Matcher Consistency):}

Let $C[j]$ be the cache entry at slot $j$, containing:
\begin{itemize}
\item $C[j].\text{sig}$: stored perturbing signature
\item $C[j].\text{valid}$: whether entry is occupied
\item $C[j].\text{decision}$: the previous decision
\end{itemize}

For an incoming operation with signature $\text{sig}_t$ and hash slot $h_t$:

If $C[h_t].\text{valid} = 1$ and $C[h_t].\text{sig} = \text{sig}_t$:
\[
\text{decision}_t = C[h_t].\text{decision}
\]
(reuse cached decision)

Otherwise:
\[
\text{decision}_t = \text{FORWARD to discrimination-detector}
\]

\textbf{Invariant:} All operations with the same perturbing signature receive the same decision.
\end{boxed}

\subsection{SystemVerilog Assertions for $P_2$}

\begin{lstlisting}[language=Verilog]
// SVA4: Cache hit correctness
property P2_cache_hit;
  @(posedge clk) 
  if (cache[hash_slot].valid && 
      cache[hash_slot].sig == incoming_sig)
    decision == cache[hash_slot].decision;
endproperty
assert property (P2_cache_hit);

// SVA5: Identical signatures reused
property P2_signature_reuse;
  @(posedge clk)
  if (sig_A == sig_B)
    (decision_A_from_cache == decision_B_from_cache);
endproperty
assert property (P2_signature_reuse);

// SVA6: Cache miss forwarding
property P2_miss_forward;
  @(posedge clk)
  if (~cache[hash_slot].valid || 
      cache[hash_slot].sig != incoming_sig)
    forward_to_detector == 1'b1;
endproperty
assert property (P2_miss_forward);

// SVA7: Cache entries persist across cycles
property P2_cache_persistence;
  @(posedge clk)
  disable iff (cache_write)
    cache_unchanged;
endproperty
assert property (P2_cache_persistence);
\end{lstlisting}

\section{Predicate $P_3$: Discrimination-Detector Refusal Tracking}

\subsection{Informal Specification}

The discrimination-detector maintains a list of perturbing signatures that have been refused in the past. When a new operation arrives (cache miss):

\begin{itemize}
\item Check if the signature is in the refusal list.
\item If yes, issue a refusal (and update the cache).
\item If no, check with the predicate controller (permission, epoch, etc.).
\end{itemize}

This enforces Axiom A2 from Doc\_0: refusal consistency.

\subsection{Formal Specification (P3)}

\begin{boxed}
\textbf{Predicate $P_3$ (Discrimination-Detector Refusal Tracking):}

Let $D$ be the set of signatures that have been refused:

\[
D = \{\text{sig} : \text{decision}(\text{sig}) = 0 \text{ in any prior cycle}\}
\]

For an incoming operation with signature $\text{sig}_t$:

If $\text{sig}_t \in D$:
\[
\text{decision}_t = 0 \text{ (refuse)}
\]

Else:
\[
\text{decision}_t = \text{consult predicate controller}
\]

\textbf{Invariant:} Once a signature is refused, it remains refused for all future operations (refusal is monotonic).
\end{boxed}

\subsection{Refusal Tracking via Registry}

The refusal set $D$ is implemented as a separate register or memory region. Updates occur only when:
\begin{itemize}
\item A new operation is evaluated by the predicate controller.
\item The predicate controller returns a refusal decision.
\item The detector adds the signature to $D$.
\end{itemize}

No removal operations are permitted.

\subsection{SystemVerilog Assertions for $P_3$}

\begin{lstlisting}[language=Verilog]
// SVA8: Previously refused signatures are refused
property P3_refusal_consistency;
  @(posedge clk)
  if (is_refused_sig(incoming_sig))
    decision == 1'b0;
endproperty
assert property (P3_refusal_consistency);

// SVA9: Refusal set monotonicity (no removals)
property P3_no_removal;
  @(posedge clk)
  if (sig_in_refused_set)
    sig_remains_in_refused_set;
endproperty
assert property (P3_no_removal);

// SVA10: Refusal update correctness
property P3_refusal_update;
  @(posedge clk)
  if (predicate_decision == REFUSE)
    add_to_refused_set(incoming_sig);
endproperty
assert property (P3_refusal_update);

// SVA11: No false positives in refusal
property P3_no_false_refuse;
  @(posedge clk)
  if (decision == REFUSE && ~is_in_refusal_set)
    predicate_controller_said_refuse;
endproperty
assert property (P3_no_false_refuse);
\end{lstlisting}

\section{Predicate $P_4$: Registry-Controller Monotonicity}

\subsection{Informal Specification}

The registry stores the perturbing set $R$. Administrative updates are received on a control interface. The controller must ensure that the perturbing set can only grow:

\[
R_t \subseteq R_{t+1} \text{ (bitwise OR)}
\]

Any attempt to remove a field is silently ignored (no error, but write has no effect).

\subsection{Formal Specification (P4)}

\begin{boxed}
\textbf{Predicate $P_4$ (Registry-Controller Monotonicity):}

The registry update rule is:

\[
R_{t+1} = R_t \mid \text{write\_request}_t
\]

where $\mid$ denotes bitwise OR. Equivalently:

\[
R_{t+1}[i] = R_t[i] \lor \text{write\_request}_t[i] \quad \forall i
\]

\textbf{Invariant:} $R_t[i] = 1 \implies R_s[i] = 1$ for all $s > t$ (once a field is marked perturbing, it remains perturbing).
\end{boxed}

\subsection{SystemVerilog Assertions for $P_4$}

\begin{lstlisting}[language=Verilog]
// SVA12: Registry update uses OR
property P4_registry_or;
  @(posedge clk)
  REG_perturb_next == (REG_perturb | reg_write_data);
endproperty
assert property (P4_registry_or);

// SVA13: Monotonicity (no removal possible)
property P4_monotonic;
  @(posedge clk)
  (REG_perturb & ~REG_perturb_next) == 0;
endproperty
assert property (P4_monotonic);

// SVA14: Once perturbing, always perturbing
property P4_persistence;
  @(posedge clk)
  if (REG_perturb[i] == 1'b1)
    REG_perturb_next[i] == 1'b1;
endproperty
assert property (P4_persistence);

// SVA15: Write-only effect is OR
property P4_write_or_effect;
  @(posedge clk)
  if (reg_write_enable == 1'b1)
    REG_perturb_next == (REG_perturb | reg_write_value);
endproperty
assert property (P4_write_or_effect);
\end{lstlisting}

\section{Cross-Predicate Invariants}

The four predicates are not independent. Their coupling enforces additional properties:

\subsection{Invariant $I_1$ (Field-Extractor → Pair-Matcher)}

The output of the field-classifier (perturbing signature) feeds directly into the pair-matcher. If $P_1$ is violated, $P_2$ cannot maintain consistency.

\subsection{Invariant $I_2$ (Pair-Matcher → Discrimination-Detector)}

Cache misses are forwarded to the discrimination-detector. If $P_2$ has a cache miss, then $P_3$ must make the decision.

\subsection{Invariant $I_3$ (Discrimination-Detector → Predicate-Combiner)}

If the discrimination-detector cannot decide (signature not in refusal set), the operation is escalated to the predicate controller. The predicate controller implements the final decision logic (permission, epoch, etc.).

\subsection{Invariant $I_4$ (Registry → Field-Classifier)}

Every time the registry $R$ is updated (via $P_4$), the field-classifier must recompute the perturbing signature using the new $R$. This ensures that $P_1$ always operates on the current perturbing set.

\section{System-Level Correctness Proof}

\begin{boxed}
\textbf{Theorem:} If $P_1 \land P_2 \land P_3 \land P_4$ hold throughout execution, then $I_0$ (master safety invariant) is preserved.

\textbf{Proof sketch:}
\begin{enumerate}
\item Assume all four predicates hold at cycle $t$.
\item Any operation $O$ at cycle $t$ has its perturbing signature computed by $P_1$.
\item If the signature is cached, $P_2$ reuses the decision (same signature $\implies$ same decision).
\item If not cached, $P_3$ checks refusal consistency (refused signature $\implies$ refuse).
\item If not refused, the predicate controller makes the decision.
\item The decision is cached for future use.
\item At any cycle $t+1$, the registry may be updated via $P_4$ (OR only).
\item All four predicates still hold at $t+1$.
\item Therefore, $I_0$ is preserved. $\square$
\end{enumerate}
\end{boxed}

\section{Formal Verification Coverage}

The 15 SVA properties (P1-P15 corresponding to SVA1-SVA15) together cover:

\begin{itemize}
\item Signature computation correctness (P1: SVA1-3)
\item Cache consistency (P2: SVA4-7)
\item Refusal tracking (P3: SVA8-11)
\item Registry monotonicity (P4: SVA12-15)
\end{itemize}

These 15 properties form a complete formal specification of the equal-status controller's behavior.

\section{Equivalence with Doc\_0 Axioms}

\begin{itemize}
\item $P_1$ enforces Axiom A1 (non-discrimination).
\item $P_3$ enforces Axiom A2 (refusal consistency).
\item $P_4$ enforces Axiom A3 (monotonicity).
\item All four predicates together enforce Axiom A4 (single-stage pipeline).
\end{itemize}

There is a bijective mapping between Doc\_0 axioms and Doc\_1 predicates. No axiom is left unimplemented. No predicate exists without an axiomatic foundation.

\section{Conclusion}

The equal-status controller is decomposed into four independent yet coupled predicates, each verified by 2-4 SVA properties. The predicates are minimal, necessary, and sufficient to preserve the master safety invariant $I_0$ from Doc\_0. All implementation details in Doc\_2 must satisfy these predicates.

\end{document}
